% Draft: §8, Conclusion, for the restructured paper (rewrite method agreed
% 17 Oct 2026; see the evidence repo's DECISIONS.md). Drafted 18 Oct from the
% structure review's plan (reviews/structure-muse-20261017-115410/MUSE.md,
% "Section 10. Conclusion": restate the limit on fusion; one paragraph on the
% companion classification; budget 300 words) and the §§1–7 drafts.
% The companion paragraph replaces the old draft's section "If determinatives
% are nouns" (fusion-of-functions-revisited.tex, sec:nouns), in the weaker form
% version A allows. Of its three points: the change of category goes (for both
% analyses, and only where a determinative heads the nominal); the secondary-use
% schema stays (CGEL pp. 520-521, 423 n. 43, checked 18 Oct); the claim that the
% pronoun and name pattern makes fusion's adjective derivation "add nothing" is
% dropped, since §4's third question now finds that pattern comes from singular
% count uses and doesn't reach plural few. The taxonomic paper
% (determinatives-as-nouns-taxonomic.tex, header comment) cites this paragraph as
% a payoff; its pointer to "§11" needs updating to §8.
%
% Label: sec:conclusion. Relies on sec:det-head, sec:mod-head, sec:once,
% sec:costs, sec:fused-relatives.
% Revised 18 Oct after the muse review (reviews/muse-sec8-20261018-124530/
% TRIAGE.md).

\section{Conclusion}\label{sec:conclusion}

On the commitments held fixed in Section~\ref{sec:analyses}, this paper adds a second function only where some property of form varies with it and not with the word. Where \textit{CGEL} fuses the head of a noun phrase with its determiner or a modifier, as in \mention{take some}, \mention{the lucky few}, and \mention{the rich}, it finds no such property, nor in Payne et al.'s \mention{once}, \mention{twice}, and \mention{thrice}, whose term-of-office uses come out as ordinary modifiers (Sections~\ref{sec:det-head}--\ref{sec:once}). On ordinary Head, \mention{few} has the same head relation in \mention{few would dare}, \mention{few of them}, \mention{the few}, and \mention{the lucky few}, where fusion gives it two different pairs of functions. Fusion's one advantage is that it derives the absence of an adjective before a bare determinative. Even there, \mention{a lucky few} and \mention{Poor everybody!} cost fusion a statement, as the condition costs ordinary Head one (Section~\ref{sec:costs}).

Noun-phrase fused relatives meet the criterion (Section~\ref{sec:fused-relatives}). The relative word's function inside its clause shows in subject relativization, in the ban on infinitival fused relatives, and, as a preference, in case. So fusion survives, favoured though not forced, as a device the grammar licenses for one construction, not as the general answer to a noun phrase with no noun.

These results hold \textit{CGEL}'s categories fixed. In separate work I argue that determinatives are a kind of noun, beside common nouns, proper nouns, and pronouns. On that classification, a nominal headed by a determinative is headed by a noun, so the change of category, Payne et al.'s second departure, goes where a determinative heads the nominal, on both analyses, and only there. It remains where an adjective heads the nominal, as in \mention{the rich}. That leaves ordinary Head with no departure from canonical structure for determinatives, while fusion keeps the first. The classification also gives a determinative the profile of a proper name: no determiner in its primary use (\mention{few would dare}), and with one, as in \mention{the few who did}, a secondary use where determiners are \enquote{selected in the ordinary way} (\mention{\bk{The Mary that you met yesterday} is my fiancée}; \citealp[520--521]{huddleston2002}). For a compound such as \mention{a special someone}, fusion would need conversion of the kind \textit{CGEL} uses for \mention{a nobody} (p.~423 n.~43), since a fused determiner-head doesn't otherwise take a determiner and an adjective before it. Nothing in the argument above depends on the classification.
